% Generated draft: requires mathematical and source review.
\documentclass[11pt]{article}
\usepackage[margin=25mm]{geometry}
\usepackage{fontspec}
\setmainfont{Latin Modern Roman}
\usepackage{amsmath,amssymb,mathtools,unicode-math}
\setmathfont{Latin Modern Math}
\usepackage{enumitem,needspace,xurl,hyperref}
\setlength{\parindent}{0pt}
\setlength{\parskip}{5pt}
\setlength{\emergencystretch}{3em}
\newcommand{\N}{\mathbb{N}}
\newcommand{\Z}{\mathbb{Z}}
\newcommand{\Q}{\mathbb{Q}}
\newcommand{\R}{\mathbb{R}}
\newcommand{\C}{\mathbb{C}}
\newcommand{\F}{\mathbb{F}}
\newcommand{\A}{\mathbb{A}}
\newcommand{\PP}{\mathbb{P}}
\newcommand{\E}{\mathbb{E}}
\newcommand{\eps}{\varepsilon}
\newcommand{\dd}{\,\mathrm{d}}
\begin{document}
\section*{UnsolvedMath ID: 3017}
\textbf{Source problem code:} catalogue entry
\subsection{Definitions and mathematical statement}
A finite presentation $\mathcal P=\langle x_1,\ldots,x_n\mid R_1,\ldots,R_m\rangle$ presents the trivial group when the normal closure of its relators in the free group $F_n$ is all of $F_n$. It is balanced when $m=n$. The source's move system is the stabilized Andrews--Curtis system:
\begin{itemize}
\item Replace $R_i$ by $R_i^{-1}$.
\item Replace $R_i$ by $R_iR_j$ for $i\ne j$.
\item Replace $R_i$ by $wR_iw^{-1}$ for a word $w$ in the generators.
\item Introduce a new generator $x_{n+1}$ and a relator $wx_{n+1}$, with $w$ a word in the existing generators, together with the inverse cancellation when applicable under the usual equivalence convention.
\end{itemize}
No arbitrary redundant relator may be added. Thus $m-n$ is unchanged. The stabilization word and inverse-move convention are made explicit as the standard meaning of a cancellable generator--relator pair.

The intended balanced stabilized conjecture asks whether every balanced presentation of the trivial group is equivalent by finitely many such transformations to a presentation $\langle x_1,\ldots,x_N\mid x_1,\ldots,x_N\rangle$ for some $N$ after stabilization. The supplied record does not state balancedness in its universal assertion, although it gives exactly that balanced target and notes invariance of $m-n$. Consequently its literal all-presentations wording needs qualification: an unbalanced presentation cannot reach that target under this move system. The balanced version above is identified openly rather than silently substituted, and it is not confused with the stronger unstabilized version that forbids generator--relator stabilization.

The source identifies this as Problem 5.2 in [Kir97]. It says the broader assertion for arbitrary finitely presented nontrivial groups is false, citing the trefoil group and [HAMS93, AC66]. It describes a $5$-dimensional handlebody built with one $0$-handle, generator $1$-handles, and relator $2$-handles; relator multiplication corresponds to a handle slide. Its asserted uniqueness and $5$-ball interpretation need balancedness and attaching/framing conventions checked, since those are not fully supplied by saying the attaching maps are homotopic. These are source-reported motivations, not a new classification of handlebodies.

The literature triage checked 2026-08-17 labels the intended trivial-group conjecture open (OPEN-TRIAGE), despite proposed counterexamples and finite searches, and reports no verified general proof or counterexample. The registry is open. No finite search or nontrivial-group counterexample is promoted to a solution of the balanced trivial-group problem.
\subsection{Short English statement}
Can every balanced trivial-group presentation be reduced to the standard trivial presentation using the source's stabilized Andrews--Curtis moves? The source reports this open, but its literal wording omits balancedness even though the moves preserve the generator--relator difference.
\subsection{Sources}
\begin{itemize}
\item Local UnsolvedMath registry, permanent ID 3017, source code catalogue entry, ``Kirby Problem 5.10''; supplied conjecture, stabilized move list, and handlebody discussion.
\item [AC65] J. J. Andrews and M. L. Curtis, ``Free groups and handlebodies,'' \emph{Proceedings of the American Mathematical Society} 16 (1965), 192--195.
\item [Kir97] R. C. Kirby, ``Problems in low-dimensional topology,'' in W. Kazez (ed.), \emph{Geometric Topology}, American Mathematical Society/International Press, Providence (1997); source identifies Problem 5.2.
\item [HAMS93] C. Hog-Angeloni, W. Metzler, and A. J. Sieradski, \emph{Two-dimensional Homotopy and Combinatorial Group Theory}, London Mathematical Society Lecture Note Series 197, Cambridge University Press (1993).
\item [AC66] J. J. Andrews and M. L. Curtis, ``Extended Nielsen operations in free groups,'' \emph{American Mathematical Monthly} 73 (1966), 21--28, doi:10.2307/2313917.
\item \emph{K3: A New Problem List in Low-Dimensional Topology}, Problem 5.10 (2026): \url{https://aimath.org/pastworkshops/kirbylistrep.pdf}.
\item Supplied literature triage checked 2026-08-17, OPEN-TRIAGE; it reports no verified resolution, records no source alteration, and cautions that dated triage does not exclude later or unindexed results.
\end{itemize}
\end{document}
